% !TeX root = ../main.tex
 
\chapter{Gruppentheorie}

\section{Körper und Ordnung auf $\bbQ$}

\dfn{Addition auf $\bbQ$}{
	Die rationalen Zahlen $\bbQ$ bilden zusammen mit der Addition
	\[
	+ \colon \bbQ \times \bbQ \to \bbQ,
	\qquad (a,b) \mapsto a+b
	\]
	eine algebraische Struktur.
}

\dfn{Gruppe}{
	Eine Menge $G$ mit einer Verknüpfung $+$ heißt \textbf{Gruppe}, falls die folgenden Eigenschaften erfüllt sind:
	\begin{enumerate}
		\item \textbf{Assoziativgesetz:}
		$\forall a,b,c \in G$ gilt
		\[
		a+(b+c)=(a+b)+c.
		\]
		\item \textbf{Neutrales Element:}
		Es existiert ein Element $e \in G$ mit
		\[
		e+a=a+e=a \qquad \forall a \in G.
		\]
		\item \textbf{Inverses Element:}
		Zu jedem $a \in G$ existiert ein Element $-a \in G$ mit
		\[
		a+(-a)=(-a)+a=e.
		\]
	\end{enumerate}
	Die Gruppe heißt \textbf{kommutativ} (oder \textbf{abelsch}), falls zusätzlich das Kommutativgesetz gilt:
	\[
	\forall a,b \in G: \quad a+b=b+a.
	\]
}

\nt{
	Die rationalen Zahlen $(\bbQ,+)$ bilden eine kommutative Gruppe mit dem neutralen Element $0$.
}

\cor{}{
	Das neutrale Element einer Gruppe ist eindeutig. Ebenso ist das inverse Element zu jedem Gruppenelement eindeutig.
}

\dfn{Multiplikation auf $\bbQ$}{
	Die rationalen Zahlen ohne die Null bilden zusammen mit der Multiplikation eine kommutative Gruppe:
	\[
	(\bbQ \setminus \{0\},\cdot).
	\]
	Das neutrale Element ist $1$, und das inverse Element zu $x \neq 0$ ist $x^{-1}=\frac{1}{x}$.
}

\dfn{Körper}{
	Ein \textbf{Körper} ist eine Menge $K$ mit zwei Verknüpfungen $+$ und $\cdot$, sodass gilt:
	\begin{enumerate}
		\item $(K,+)$ ist eine kommutative Gruppe mit neutralem Element $0$.
		\item $(K \setminus \{0\},\cdot)$ ist eine kommutative Gruppe mit neutralem Element $1$.
		\item Es gilt das \textbf{Distributivgesetz}:
		\[
		\forall x,y,z \in K:
		\quad x\cdot(y+z)=x\cdot y+x\cdot z.
		\]
	\end{enumerate}
}

\nt{
	Die rationalen Zahlen $\bbQ$ bilden mit der Addition und der Multiplikation einen Körper.
}

Aus den Körperaxiomen lassen sich elementare Rechenregeln herleiten.

\mlemma{Rechenregeln}{
	Für alle $x,y \in K$ gilt:
	\[
	x\cdot 0=0,
	\qquad
	(-x)\cdot y=-(x\cdot y).
	\]
}

\begin{myproof}
	Mit dem Distributivgesetz folgt
	\begin{align*}
		x\cdot 0
		&=x\cdot(0+0)\\
		&=x\cdot 0+x\cdot 0.
	\end{align*}
	Durch Addition des inversen Elements $-(x\cdot 0)$ auf beiden Seiten erhalten wir
	\[
	x\cdot 0=0.
	\]
	Weiter gilt
	\begin{align*}
		x\cdot y+(-x)\cdot y
		&=(x+(-x))\cdot y\\
		&=0\cdot y\\
		&=0.
	\end{align*}
	Somit ist $(-x)\cdot y$ das additive Inverse von $x\cdot y$. Wegen der Eindeutigkeit des inversen Elements folgt
	\[
	(-x)\cdot y=-(x\cdot y).
	\]
\end{myproof}

\section{Ordnung auf $\bbQ$}

\dfn{Ordnung}{
	Eine Relation $<$ auf einem Körper $K$ heißt \textbf{vollständige Ordnung}, falls für alle $x,y,z \in K$ gilt:
	\begin{enumerate}
		\item \textbf{Trichotomie:} Genau eine der drei Aussagen
		\[
		x<y,\qquad x=y,\qquad x>y
		\]
		ist erfüllt.
		\item \textbf{Transitivität:}
		\[
		x<y \land y<z \implies x<z.
		\]
		\item \textbf{Kompatibilität mit der Addition:}
		\[
		x<y \implies x+z<y+z.
		\]
		\item \textbf{Kompatibilität mit der Multiplikation:}
		\[
		x<y \land z>0 \implies xz<yz.
		\]
	\end{enumerate}
}

Der Körper $\bbQ$ besitzt eine vollständige Ordnung, die mit der Addition und der Multiplikation verträglich ist.

\nt{
	Für alle $x,y \in \bbQ$ gilt
	\[
	x<y \iff 0<y-x.
	\]
}

\mlemma{Positive rationale Zahlen}{
	Für alle $x,y \in \bbQ$ gilt:
	\begin{enumerate}
		\item Sind $x>0$ und $y>0$, so ist $x+y>0$.
		\item Sind $x>0$ und $y>0$, so ist $xy>0$.
	\end{enumerate}
}

\cor{Transitivität der Ordnung}{
	Aus $x<y$ und $y<z$ folgt $x<z$.

	\begin{myproof}
		Aus den Voraussetzungen folgt
		\[
		0<y-x,\qquad 0<z-y.
		\]
		Durch Addition erhalten wir
		\[
		0<(y-x)+(z-y)=z-x.
		\]
		Also gilt $x<z$.
	\end{myproof}
}

\cor{Multiplikation mit positiven Zahlen}{
	Für $x,y,z \in \bbQ$ gilt
	\[
	x<y \land z>0 \implies xz<yz.
	\]
	\begin{myproof}
		Aus $x<y$ folgt $0<y-x$. Wegen $z>0$ gilt
		\[
		0<z(y-x)=zy-zx.
		\]
		Damit folgt $xz<yz$.
	\end{myproof}
}

\cor{Quadrate sind nichtnegativ}{
	Für alle $x\in\bbQ$ gilt $x^2\geq 0$. Insbesondere gilt für $x\neq 0$:
	\[
	x^2>0.
	\]
	\begin{myproof}
		Wir unterscheiden zwei Fälle:
		\begin{enumerate}
			\item Falls $x>0$, folgt aus der Verträglichkeit der Ordnung mit der Multiplikation $x^2>0$.
			\item Falls $x<0$, gilt $-x>0$. Somit ist
			\[
			x^2=(-x)(-x)>0.
			\]
		\end{enumerate}
		Für $x=0$ gilt $x^2=0$. Insgesamt folgt $x^2\geq 0$ für alle $x\in\bbQ$.
	\end{myproof}
}

\cor{}{
	Insbesondere gilt
	\[
	1>0,
	\]
	da $1=1^2$ und $1\neq 0$.
}

\cor{Es gibt stets eine größere rationale Zahl}{
	Für jedes $q\in\bbQ$ existiert ein $q'\in\bbQ$ mit
	\[
	q<q'.
	\]

	\begin{myproof}
		Zunächst gilt $0<1$. Aufgrund der Kompatibilität der Ordnung mit der Addition folgt
		\[
		q=q+0<q+1.
		\]
		Wir wählen also $q'=q+1\in\bbQ$. 
	\end{myproof}
}

\cor{Zwischen $0$ und jeder positiven rationalen Zahl liegt eine weitere rationale Zahl}{
	Für jedes $q\in\bbQ$ mit $q>0$ existiert ein $q'\in\bbQ$ mit
	\[
	0<q'<q.
	\]

	\begin{myproof}
		Sei $q>0$. Dann gilt auch $q+1>0$, sodass $(q+1)^{-1}>0$ existiert.

		Da $q<q+1$ und $(q+1)^{-1}>0$, folgt aus der Kompatibilität der Ordnung mit der Multiplikation:
		\[
		\frac{q}{q+1}<1.
		\]
		Insbesondere gilt
		\[
		0<\frac{q}{q+1}<1.
		\]
		Wir setzen
		\[
		q'=\frac{q}{q+1}\cdot q=\frac{q^2}{q+1}.
		\]
		Dann ist $q'>0$. Außerdem gilt wegen $q<q+1$:
		\[
		\frac{q}{q+1}<1
		\implies
		\frac{q^2}{q+1}<q.
		\]
		Somit erhalten wir
		\[
		0<q'<q.
		\]
	\end{myproof}
}

\dfn{Absolutbetrag}{
	$\vert x \vert =
	\begin{cases}
		\vert x \vert &= x \text{ falls } x \ge 0 \\
		\vert x \vert &= -x \text{ falls } x < 0
	\end{cases}$
}

\mlemma{}{
	\begin{enumerate}
		\item $\vert x \vert = 0 \Leftrightarrow x = 0$
		\item $\vert x \cdot y \vert = \vert x \vert \cdot \vert y \vert$.
		\item \textbf{Dreiecksungleichung} $\vert x + y \vert \leq \vert x \vert + \vert y \vert.$
		\item $\vert \vert x \vert - \vert y \vert \vert \leq \vert x + y \vert $
	\end{enumerate}
}
\begin{myproof}
	Wegen $\pm x \leq \vert x \vert \land \pm y \leq \vert y \vert$ gilt auch $\pm \vert (x + y) \vert$ 
\end{myproof}

\newpage
\section{Lösungen von Gleichungen}

\qs{}{
	Beweise, dass die Gleichung $x^2=2$ keine Lösung in $\bbQ$ besitzt.
}

\begin{myproof}
	Wir führen einen Widerspruchsbeweis.

	Angenommen, es existiert eine rationale Zahl $x\in\bbQ$ mit
	\[
	x^2=2.
	\]
	Dann lässt sich $x$ als vollständig gekürzter Bruch
	\[
	x=\frac{p}{q}, \qquad p,q\in\bbN,\quad \gcd(p,q)=1
	\]
	schreiben. Durch Quadrieren erhalten wir
	\begin{align*}
		\frac{p^2}{q^2}&=2\\
		\implies p^2&=2q^2.
	\end{align*}
	Somit ist $p^2$ gerade. Daher ist auch $p$ gerade, also existiert ein $k\in\bbN$ mit $p=2k$.

	Einsetzen liefert
	\begin{align*}
		(2k)^2&=2q^2\\
		4k^2&=2q^2\\
		\implies q^2&=2k^2.
	\end{align*}
	Also ist auch $q^2$ gerade und damit $q$ gerade.

	Dies bedeutet, dass sowohl $p$ als auch $q$ durch $2$ teilbar sind. Das widerspricht der Annahme $\gcd(p,q)=1$.

	Somit besitzt die Gleichung $x^2=2$ keine rationale Lösung.
\end{myproof}

\section{Supremum und Vollständigkeit}

\dfn{Obere Schranke und Supremum}{
	Sei $A\subseteq\bbR$ eine nichtleere Menge.

	\begin{itemize}
		\item Die Menge $A$ heißt \textbf{nach oben beschränkt}, falls es ein $s\in\bbR$ gibt, sodass
		\[
		\forall x\in A:\quad x\leq s.
		\]
		\item Ein solches $s$ heißt \textbf{obere Schranke} von $A$.
		\item Eine obere Schranke $s_0$ heißt \textbf{Supremum} von $A$, falls sie kleiner oder gleich jeder anderen oberen Schranke ist:
		\[
		\forall s\in\bbR:
		\quad
		s\text{ obere Schranke von }A
		\implies s_0\leq s.
		\]
		Wir schreiben
		\[
		s_0=\sup A.
		\]
	\end{itemize}
}

\ex{Supremum}{
	Wir betrachten die Mengen
	\[
	A=[0,1]=\{x\in\bbR\mid 0\leq x\leq 1\}
	\]
	und
	\[
	B=[0,1)=\{x\in\bbR\mid 0\leq x<1\}.
	\]

	Für beide Mengen gilt
	\[
	\sup A=\sup B=1.
	\]

	Bei $A$ ist $1$ sogar das größte Element der Menge, da $1\in A$ gilt.

	Bei $B$ ist $1$ dagegen kein Element der Menge. Dennoch ist $1$ die kleinste obere Schranke.

	\begin{myproof}
		Wir zeigen, dass $1=\sup B$ gilt.

		Zunächst ist $1$ eine obere Schranke von $B$, denn für alle $x\in B$ gilt $x<1$.

		Sei nun $s<1$ eine beliebige Zahl. Wir unterscheiden zwei Fälle:
		\begin{enumerate}
			\item Falls $s<0$, ist $s$ keine obere Schranke von $B$, da $0\in B$ gilt.
			\item Falls $0\leq s<1$, setzen wir
			\[
			\varepsilon=\frac{1-s}{2}>0.
			\]
			Dann gilt
			\[
			s<s+\varepsilon=\frac{1+s}{2}<1.
			\]
			Somit ist $s+\varepsilon\in B$ und $s$ keine obere Schranke von $B$.
		\end{enumerate}
		Also ist keine Zahl $s<1$ eine obere Schranke von $B$. Folglich ist $1$ die kleinste obere Schranke und damit
		\[
		\sup B=1.
		\]
	\end{myproof}
}

\dfn{Vollständigkeitsaxiom für $\bbR$}{
	Die reellen Zahlen $\bbR$ heißen \textbf{vollständig}, falls jede nichtleere, nach oben beschränkte Teilmenge $A\subseteq\bbR$ ein Supremum besitzt, das heißt:
	\[
	\forall A\subseteq\bbR:
	\quad
	\left.
	\begin{array}{l}
		A\neq\emptyset,\\
		A\text{ ist nach oben beschränkt}
	\end{array}
	\right\}
	\implies \sup A\in\bbR.
	\]
}
